\documentclass[12pt,a4paper]{article}

% ---- Język i kodowanie ----
\usepackage{iftex}
\ifPDFTeX
  \usepackage[utf8]{inputenc}
  \usepackage[T1]{fontenc}
  \usepackage{lmodern}
\else
  \usepackage{fontspec}
\fi
\usepackage[polish]{babel}
\usepackage{indentfirst}

% ---- Układ strony ----
\usepackage{geometry}
\geometry{margin=2.5cm}
\usepackage{microtype}
\setlength{\parskip}{0.4em}

% ---- Tabele, listy ----
\usepackage{float}
\usepackage{caption}
\captionsetup{font=small,labelfont=bf}
\usepackage{booktabs}
\usepackage{tabularx}
\usepackage{longtable}
\usepackage{array}
\usepackage{enumitem}
\setlist{nosep,leftmargin=1.4em}
\newcolumntype{Y}{>{\raggedright\arraybackslash}X}

% ---- Kolory (spójne z raportem końcowym PlayLink) ----
\usepackage{xcolor}
\definecolor{gold}{HTML}{8A6D3B}
\definecolor{darkgold}{HTML}{5A4A2A}
\newcommand{\code}[1]{\texttt{\detokenize{#1}}}

% ---- Diagram ----
\usepackage{tikz}
\usetikzlibrary{positioning,arrows.meta,fit,backgrounds}

% ---- Tytuły sekcji ----
\usepackage{titlesec}
\titleformat{\section}{\Large\bfseries\color{darkgold}}{\thesection}{0.6em}{}
\titleformat{\subsection}{\large\bfseries\color{gold}}{\thesubsection}{0.6em}{}

% ---- Hiperłącza ----
\usepackage[hidelinks]{hyperref}
\hypersetup{colorlinks=true,linkcolor=darkgold,urlcolor=gold}

\title{
    \large Projekt kompetencyjny IOAD 2026/2027 \\[0.3em]
    \Huge \textbf{Dokument wizji} \\[0.3em]
    \LARGE PlayLink Multi-Region \\
    \large \textit{Lobby w dwóch regionach chmury: routing, powtarzalne wdrożenie, eksperyment i próby awarii}
}
\author{
    Bartosz Kołaciński, nr indeksu: 251554 \\
    Mateusz Kosowski, nr indeksu: 251558 \\
    Nikodem Nowak, nr indeksu: 251598 \\
    Natan Ręklewski, nr indeksu: 252952 \\
    Jakub Rosiak, nr indeksu: 251620 \\
    Andrzej Krysiak, nr indeksu: 252945 \\[0.6em]
    \normalsize Opiekun projektu: mgr inż. Bartłomiej Jencz
}
\date{\today}

\begin{document}
\pagenumbering{gobble}
\maketitle
\vfill
\begin{center}
\small
Repozytorium: \url{https://github.com/ioad-modul-sumatywny-26/playlink}
\end{center}

\newpage
\pagenumbering{Roman}
\tableofcontents
\newpage

% =====================================================================
\pagenumbering{arabic}
\section{Informacje o dokumencie}

\subsection{Cel dokumentu}
Dokument przedstawia wizję rozwoju aplikacji PlayLink z pojedynczej instancji do
platformy działającej w dwóch regionach chmury publicznej (docelowo dwa oddalone od siebie regiony, np. Europa i USA). Stos
regionalny stawiamy na Microsoft Azure (kredyty studenckie); AWS jest alternatywą.
Domena, tunele i ochrona przed atakami są na Cloudflare. Maszyny i kontenery,
razem z bazą, są na Azure. Wybór regionu robi przeglądarka,
bez trzeciego serwera (sekcja~\ref{sec:chmura}). Dokument opisuje
problem biznesowy, wartość dla użytkowników, zakres, wymagania, ryzyka i kryteria
sukcesu, uwzględniając zrealizowane rezultaty prac (PR 91, 92, 96, 97, 99): wdrożenie wieloregionowe EU/US za tunelami Cloudflare, routing z propozycją Stay/Switch w przeglądarce, zautomatyzowane wdrożenie pull-based z GHCR, eksperyment obciążeniowy F4, centralny stos obserwowalności (Grafana, VictoriaLogs, VictoriaMetrics) oraz dedykowany serwer Model Context Protocol (MCP).

\subsection{Definicje i skróty}
\begin{table}[H]
\centering\small
\begin{tabularx}{\textwidth}{@{}lY@{}}
\toprule
\textbf{Termin} & \textbf{Znaczenie} \\
\midrule
LFG & \emph{Looking For Group} --- szukanie współgraczy. \\
Lobby / pokój & Krótkotrwała sesja gry z limitem miejsc, wygasająca po czasie. \\
Region & Osobne środowisko w chmurze (np. Frankfurt, N.\ Virginia)  z kompletnym stosem aplikacji. \\
Stos regionalny & Frontend, backend i baza danych wdrożone w jednym regionie. \\
Routing użytkowników & Mechanizm wyboru regionu dla użytkownika (automatyczny lub ręczny). \\
Stay/Switch & Mechanizm propozycji zmiany regionu: jeśli zapisany region żyje, interfejs proponuje szybszy region banerem bez agresywnego przerywania sesji. \\
IaC & \emph{Infrastructure as Code} --- opis maszyn, bazy, DNS i sekretów w plikach w repozytorium. Jedno polecenie tworzy region; kolejny region to nowy plik parametrów, a nie klikanie w panelu. \\
Sonda regionu & Pomiar czasu odpowiedzi endpointu zdrowia, wykonywany w przeglądarce przed wyborem regionu. \\
RTT & \emph{Round-Trip Time} --- czas przejścia sygnału tam i z powrotem. \\
RTO / RPO & Maksymalny akceptowalny czas niedostępności / utraty danych po awarii. \\
p50, p95 & Mediana i 95.\ percentyl czasu odpowiedzi. \\
WS & WebSocket --- kanał czasu rzeczywistego (lista pokojów, czat). \\
MCP & \emph{Model Context Protocol} --- standard bezpiecznej komunikacji narzędziowej i udostępniania zasobów dla asystentów AI (Claude, Antigravity, Cursor). \\
VictoriaLogs / Metrics & Centralny system składowania i przeszukiwania logów (LogsQL) oraz metryk (PromQL) ze wszystkich regionów. \\
\bottomrule
\end{tabularx}
\end{table}

% =====================================================================
\newpage
\section{Problem biznesowy}

\subsection{Obecna sytuacja}
\textbf{Jak działa dziś.} PlayLink to działająca aplikacja LFG (FastAPI, SvelteKit,
PostgreSQL, Docker Compose) wdrażana automatycznie z gałęzi \code{main} na
\emph{jeden} serwer. Każdy pokój ma pole \code{lobby_location}, ale to tylko
\emph{etykieta} (10 predefiniowanych lokalizacji w \code{backend/config.py}); nie
oznacza ona, że pokój obsługuje serwer w tym regionie. Wskaźnik ping w interfejsie
mierzy czas do tego jednego backendu, więc nie mówi nic o jakości połączenia z
lobby w innym regionie.

\textbf{Jakie są problemy.}
\begin{itemize}
  \item Użytkownik z daleka od serwera ma wyższe opóźnienia przy każdej akcji
        (utworzenie pokoju, dołączenie, czat, odświeżenie listy).
  \item Lobby „w USA'' i „w Europie'' są rozróżnione tylko opisowo --- nie da się
        zweryfikować, czy dwa regiony faktycznie pomagają graczom.
  \item Pojedynczy punkt awarii: niedostępność serwera wyłącza całą platformę.
  \item Stan czasu rzeczywistego jest trzymany w pamięci procesu
        (menedżery połączeń WebSocket i liczniki limitów
        czatu), więc obecna architektura nie skaluje się poziomo ``za darmo''.
  \item Wdrożenie jest opisane skryptem CI dla jednego hosta, a nie powtarzalnym
        szablonem infrastruktury.
\end{itemize}

\textbf{Koszty i ryzyka.} Wyższe opóźnienia zniechęcają graczy z odległych
regionów (a w grach niszowych każdy odpadający gracz osłabia i tak małą
społeczność). Awaria jedynego serwera oznacza całkowity brak usługi. Rozbudowa „na
oko'' o drugi region bez pomiarów mogłaby podwoić koszty bez udowodnionej korzyści.

\subsection{Problem do rozwiązania}
Nie wiemy, \emph{czy i kiedy} wdrożenie w dwóch regionach realnie poprawia
doświadczenie graczy, jaką cenę (koszt, złożoność, ograniczenia) trzeba za to
zapłacić i jak zachowa się system, gdy jeden region przestanie działać. Projekt ma
na te pytania odpowiedzieć \textbf{pomiarem}, a nie założeniem.

% =====================================================================
\section{Wizja produktu}
\begin{quote}
Dla \textbf{graczy gier niszowych i retro z Europy i Ameryki Północnej}, którzy
potrzebują \textbf{szybkiego, niezawodnego lobby blisko siebie}, produkt
\textbf{PlayLink Multi-Region} jest \textbf{platformą LFG działającą w dwóch
regionach chmury z automatycznym i ręcznym wyborem regionu}, która umożliwia
\textbf{niskie opóźnienia, ciągłość działania przy awarii jednego regionu oraz
oparte na pomiarach decyzje o tym, czy dodatkowy region jest opłacalny}.
\end{quote}

\subsection{Jedna idea, sześć odrębnych rezultatów}
Wszystkie elementy poniżej służą jednej idei (dwa regiony), ale każdy odpowiada na
\emph{inne pytanie}, ma \emph{inny produkt końcowy} i \emph{inną miarę sukcesu}.
Dzięki temu nie są sześcioma opisami tej samej rzeczy.

\begin{table}[H]
\centering\footnotesize
\renewcommand{\arraystretch}{1.3}
\begin{tabularx}{\textwidth}{@{}l>{\raggedright\arraybackslash}p{0.2\textwidth}Y>{\raggedright\arraybackslash}p{0.2\textwidth}@{}}
\toprule
\textbf{Filar} & \textbf{Pytanie} & \textbf{Co zrealizowano} & \textbf{Miara} \\
\midrule
F1 Regiony i lobby & Co dokładnie działa w regionie? & Dwa kompletne stosy (EU, US) na Azure z własnymi lobby za Cloudflare Tunnel & Oba regiony obsługują pełny przepływ: pokój, dołączenie, czat (PR 91) \\
F2 Routing & Skąd system wie, dokąd wysłać użytkownika? & Sonda opóźnień w przeglądarce, baner propozycji Stay/Switch i ręczny przełącznik & Użytkownik trafia do szybszego regionu; brak niechcianych przeskoczeń (PR 91, 92) \\
F3 Wdrożenie & Czy środowisko da się odtworzyć bez ręcznej pracy? & Szablony compose, skrypty \code{deploy/deploy.sh} i wdrożenie pull-based z GHCR na Azure VM & Czysty start regionu w kilka minut; automatyczna aktualizacja wydań (PR 91) \\
F4 Eksperyment & Czy dwa regiony są szybsze niż jeden? & Dedykowany runner i scenariusze w \code{experiments/f4/}; zebrane pomiary 1 vs 2 regionów & Zmierzona różnica p50/p95, błędów i czasu synchronizacji WebSocket (PR 92) \\
F5 Awarie & Co się dzieje, gdy region pada? & Scenariusze awarii, procedury runbookowe, automatyczne przełączenie przy braku odpowiedzi & Czas wykrycia, ciągłość działania drugiego regionu i brak wpływu na bazę \\
F6 Analiza & Kiedy to się opłaca? & Raport z wnioskami: korzyści, koszty maszyn w Azure, transfer i podział graczy & Rekomendacja poparta rzeczywistym rachunkiem i pomiarami F4 \\
\bottomrule
\end{tabularx}
\caption{Odrębność sześciu filarów projektu oraz ich stan realizacji.}
\end{table}

\subsection{Zależności między filarami}
F1 jest fundamentem (bez regionów nie ma czego routować). F2 i F3 są niezależnymi
strumieniami pracy nad F1 (kod aplikacji vs.\ infrastruktura). F4 wymaga F1--F3 i
produkuje dane dla F6. F5 wymaga F1--F3, a jej wyniki są osobnym wejściem do F6.
Dzięki temu zespół może pracować równolegle, a F4--F6 są etapem \emph{badawczym},
nie kolejnym etapem budowy.

% =====================================================================
\section{Cele biznesowe}
\begin{table}[H]
\centering\small
\begin{tabularx}{\textwidth}{@{}lYY@{}}
\toprule
\textbf{ID} & \textbf{Cel} & \textbf{Sposób weryfikacji} \\
\midrule
CB-1 & Skrócić opóźnienia dla graczy z obu kontynentów w porównaniu z jednym regionem & Eksperyment F4: p95 dla użytkowników z USA i z Europy \\
CB-2 & Zachować możliwość gry mimo awarii jednego regionu & Próby awarii F5: czas odzyskania, dostępność drugiego regionu \\
CB-3 & Umożliwić odtworzenie środowiska jednym poleceniem & F3: czysty start w nowym koncie lub regionie \\
CB-4 & Poznać koszt i ograniczenia modelu dwuregionowego & F6: tabela kosztów i ograniczeń \\
CB-5 & Utrzymać dotychczasowe wymagania jakościowe (WN-01 p95 $<$ 200~ms, WN-02 $<$ 1000~ms) & Regresja testów i pomiary w obu regionach \\
\bottomrule
\end{tabularx}
\end{table}

% =====================================================================
\section{Grupy użytkowników}
\begin{table}[H]
\centering\small
\begin{tabularx}{\textwidth}{@{}lYY@{}}
\toprule
\textbf{Grupa} & \textbf{Charakterystyka} & \textbf{Potrzeba w wizji} \\
\midrule
Gracz z Europy & Najliczniejsza grupa dziś; blisko obecnego serwera & Brak pogorszenia jakości; możliwość dołączenia do lobby w USA \\
Gracz z USA & Dziś obsługiwany z daleka & Lobby blisko siebie i niższe opóźnienia \\
Gracz „międzyregionalny'' & Gra w tytuły o małej społeczności, chce większej puli graczy & Ręczny wybór regionu niezależnie od lokalizacji \\
Administrator & Moderuje pokoje i kategorie gier & Zarządzanie w każdym regionie \\
Operator / zespół & Wdraża i obserwuje system & Jeden szablon, metryki, alarmy \\
\bottomrule
\end{tabularx}
\end{table}

% =====================================================================
\section{Zakres produktu}

\subsection{Funkcjonalności w zakresie}

\paragraph{F1 --- Dwa regiony chmury z działającymi lobby.}
Zrealizowano dwa niezależne stosy regionalne (Europa i USA, projekty \code{playlink-eu} i \code{playlink-us} w PR 91), każdy z frontendem, backendem i
bazą PostgreSQL. \textbf{Lobby jest regionalne}: pokój istnieje w bazie regionu, w
którym go utworzono. To świadoma decyzja architektoniczna --- stan czasu
rzeczywistego jest w pamięci procesu, a pokoje są krótkotrwałe, więc
replikacja między regionami dawałaby dużą złożoność za niewielki zysk. Maszyny wirtualne nie posiadają otwartych portów publicznych --- ruch w całości obsługuje Cloudflare Tunnel. Endpointy \code{/healthz} i \code{/readyz} weryfikują stan backendu i bazy, a rate limiter opiera się na nagłówku \code{CF-Connecting-IP}.
\emph{Wspólna} jest tożsamość (konto powstaje przy pierwszym podpisie, a klucz nie
opuszcza przeglądarki), \emph{regionalne} są pokoje, czat i listy.

\newpage
\paragraph{F2 --- Kierowanie użytkowników.}
Adres strony wskazuje na jeden z regionów za pośrednictwem Cloudflare. Skrypt w odpowiedzi
mierzy czas do endpointu zdrowia każdego regionu i otwiera
lobby w szybszym, który odpowiedział. Osobnego serwera routingu
nie ma. Wybór regionu jest
zapamiętywany w cookie (\code{playlink.region}). W ramach PR 92 wdrożono dopracowany mechanizm \textbf{Stay/Switch} (komponent \code{RegionProposalBanner}): przy kolejnych wejściach gracz wraca do zapisanego regionu, a gdy pomiar w tle wskaże, że drugi region jest szybszy, interfejs wyświetla nienatarczywy baner z propozycją zmiany. Sesja nie jest przerywana samoczynnie. Dopiero gdy zapisany region nie odpowiada, następuje automatyczny failover na żywy region. Ręczny przełącznik w nagłówku pozwala na swobodną zmianę w dowolnym momencie. Adresy API i WS są dynamicznie konfigurowane pod wybrany region.

\paragraph{F3 --- Powtarzalne wdrożenie i automatyzacja wydań.}
Zamiast ręcznego wdrożenia czy bezpośredniego dostępu SSH z potoku CI, wdrożono zautomatyzowany, bezpieczny model pull-based (PR 91):
\begin{itemize}
  \item Skrypt \code{deploy/deploy.sh <eu|us> [tag]} umożliwia natychmiastowe uruchomienie parametryzowanego stosu regionalnego (lokalnie lub w chmurze) z wygenerowanym środowiskiem \code{.run/<region>.env}.
  \item Potok wydań \code{.github/workflows/release.yml} buduje wieloarchitekturowe obrazy i publikuje je w GitHub Packages (GHCR) po oznaczeniu commita tagiem wydania.
  \item Maszyny Azure mają skonfigurowany timer systemd (\code{deploy/update.sh} co 2 minuty), który automatycznie pobiera najnowszy release (najpierw EU, po weryfikacji zdrowia US). W razie błędu healthchecku skrypt automatycznie wycofuje wdrożenie (rollback).
  \item Dla wygody prac deweloperskich przygotowano lokalne środowisko wieloregionowe (\code{start-local.ps1} z PR 92).
\end{itemize}

\paragraph{F4 --- Eksperyment porównawczy.}
W ramach PR 92 zrealizowano dedykowany moduł eksperymentu obciążeniowego w katalogu \code{experiments/f4/}. Przygotowano syntetycznego klienta asynchronicznego oraz skrypty runnera (\code{run.py}, \code{run-f4.ps1}), które porównują wariant A (pojedynczy region EU) z wariantem B (dwa regiony EU/US). Scenariusze obejmują tworzenie pokojów, dołączanie, czat WebSocket oraz propagację listy. Narzędzia mierzą opóźnienia (p50, p95), kody odpowiedzi, błędy i zachowanie pod obciążeniem, generując ustrukturyzowane raporty JSON i Markdown.

\paragraph{F5 --- Próby awarii.}
Kontrolowane wyłączenie jednego regionu (zatrzymanie usługi backendu, odcięcie
ruchu tunelu, zatrzymanie bazy), obserwacja drugiego i pomiar odzyskania: czas wykrycia,
czas przełączenia ruchu, zachowanie klientów (frontend posiada heartbeat WS 30~s i
automatyczne ponowne łączenie) oraz stan pokoi po powrocie regionu (udokumentowane w \code{docs/operations/runbook-regions.md}).

\paragraph{F6 --- Analiza wyników.}
Raport z wnioskami: kiedy dwa regiony pomagają, ile kosztują (obliczenia, baza,
transfer międzyregionalny, sonda regionu) i jakie mają ograniczenia (w tym
podział społeczności graczy między regiony), zasilany danymi z eksperymentu F4 i prób F5.

\paragraph{F7 --- Obserwowalność i inspekcja MCP.}
W ramach PR 96, 97 i 99 rozszerzono zakres platformy o zaawansowane narzędzia diagnostyczne i telemetryczne:
\begin{itemize}
  \item \textbf{Centralna obserwowalność (PR 99):} Uruchomiono maszynę \code{playlink-ops} ze stosem VictoriaLogs (logi), VictoriaMetrics (metryki) i Grafana (\code{deploy/observability.sh}). Kontenery Vector w każdym regionie przesyłają logi JSON z identyfikatorem \code{request_id} i zdarzeniami biznesowymi, a cAdvisor zbiera metryki zużycia zasobów. Błędy frontendu trafiają do logów przez \code{/api/client-errors}. Panel Grafany chroniony jest przez Cloudflare Access.
  \item \textbf{Serwer MCP (PR 96, 97, 99):} Wdrożono bezpieczny serwer MCP w trybie \emph{read-only} (\code{backend/mcp_server/}) z transportem stdio i SSE. Oferuje narzędzia inspekcji gier, pokojów, zdrowia (\code{check_platform_health}) oraz zasoby per region (\code{playlink://{region}/*}) z agregacją multi-region. Dostępne są też narzędzia kwerend logów LogsQL (\code{search_logs}) oraz metryk PromQL (\code{query_metrics}).
\end{itemize}

\subsection{Poza zakresem}
\begin{itemize}
  \item Replikacja bazy i pokoi między regionami (aktywna--aktywna) oraz
        współdzielone lobby ponad regionami.
  \item Uruchomienie trzeciego regionu w tych 10 tygodniach. Szablon i lista
        adresów mają go przyjąć bez zmiany architektury; samo wdrożenie zostaje
        poza zakresem, o ile po F1--F6 nie zostanie czas.
  \item Przeniesienie aplikacji na Cloudflare Containers albo bazy na D1
        (uzasadnienie w sekcji~\ref{sec:chmura}).
  \item Osobna usługa balansowania (Azure Traffic Manager, Cloudflare Load
        Balancing) na ścieżce ruchu. Zostają opisane jako odpowiednik produkcyjny
        sondy w przeglądarce.
  \item Autoskalowanie produkcyjne, CDN poza podstawowym, WAF i pełna zgodność
        z regulacjami (np.\ RODO dla danych w USA) --- odnotowane jako ryzyko.
  \item Matchmaking oparty na umiejętnościach, płatności, aplikacje mobilne.
  \item Zmiana mechanizmu uwierzytelniania (BIP39/ECDSA pozostaje bez zmian).
\end{itemize}

% =====================================================================
\section{Kluczowe wymagania biznesowe}
\begin{table}[H]
\centering\small
\begin{tabularx}{\textwidth}{@{}lYl@{}}
\toprule
\textbf{ID} & \textbf{Wymaganie} & \textbf{Filar} \\
\midrule
WB-1 & Każdy region obsługuje pełny przepływ: logowanie, tworzenie, dołączanie, czat, lista na żywo & F1 (PR 91) \\
WB-2 & Przy pierwszym wejściu użytkownik trafia do szybszego żywego regionu, zanim otworzy się lobby & F2 (PR 91) \\
WB-3 & Zapisany region zostaje. Szybszy drugi region jest propozycją (Stay/Switch). Awaria zapisanego regionu przełącza na drugi. Ręczna zmiana działa & F2 (PR 91, 92) \\
WB-4 & Oba środowiska uruchamia jeden szablon lub skrypt, z automatycznym wdrożeniem wydań pull-based & F3 (PR 91) \\
WB-5 & Porównanie 1 vs 2 regionów odbywa się według opisanego, powtarzalnego scenariusza w \code{experiments/f4/} & F4 (PR 92) \\
WB-6 & Awaria jednego regionu nie wyłącza drugiego, a procedury odzyskania są udokumentowane w runbooku & F5 \\
WB-7 & Wnioski zawierają koszt i ograniczenia, nie tylko korzyści & F6 \\
WB-8 & Konto użytkownika działa w obu regionach bez ponownej rejestracji & F1/F2 \\
WB-9 & Centralna obserwowalność (logi, metryki, Grafana) oraz inspekcja platformy przez asystentów AI (MCP) & F7 (PR 96, 97, 99) \\
\bottomrule
\end{tabularx}
\end{table}

% =====================================================================
\section{Główne procesy biznesowe}

\begin{enumerate}[label=\textbf{P\arabic*.},leftmargin=2.2em]
  \item \textbf{Wejście użytkownika.} Adres strony otwiera jeden z regionów
        i stamtąd przychodzi strona ze skryptem sondy. Skrypt mierzy ping i dopiero
        potem otwiera lobby w szybszym żywym regionie. Kolejna wizyta wraca do
        zapisanego regionu. Logowanie to podpis tym samym kluczem.
        Pierwszy podpis w regionie zakłada tam konto z losową nazwą.
  \item \textbf{Zmiana regionu.} Gracz zmienia region ręcznie albo przyjmuje
        propozycję Stay/Switch, gdy drugi region jest szybszy. Adresy API i WS przełączają się,
        a lista pokojów jest wczytywana z nowego regionu. Sam pomiar regionu nie zmienia bez zgody użytkownika.
  \item \textbf{Gra w lobby.} Tworzenie i dołączanie do pokoju, czat i RSVP
        działają w obrębie wybranego regionu; pokoje wygasają jak dotychczas.
  \item \textbf{Wdrożenie i aktualizacja.} Publikacja wydań na GHCR uruchamia pull-based timer systemd na maszynach Azure (\code{deploy/update.sh}), który pobiera obrazy i bezpiecznie aktualizuje regiony (najpierw EU, po weryfikacji zdrowia US).
  \item \textbf{Eksperyment.} Uruchomienie scenariuszy obciążeniowych z katalogu \code{experiments/f4/} dla wariantu 1 i 2 regionów; rejestracja czasów p50/p95, błędów i synchronizacji WebSocket.
  \item \textbf{Awaria i odzyskanie.} Wyłączenie regionu, automatyczne przekierowanie graczy na zdrowy region i weryfikacja procedur awaryjnych z runbooka.
  \item \textbf{Analiza.} Zestawienie metryk, kosztów maszyn w Azure oraz wniosków architektonicznych.
  \item \textbf{Obserwowalność i inspekcja MCP.} Zbieranie logów i metryk przez Vector/cAdvisor do VictoriaLogs/VictoriaMetrics oraz inspekcja stanu i telemetrii przez serwer MCP.
\end{enumerate}

\begin{figure}[H]
\centering
\begin{tikzpicture}[
  font=\small,
  box/.style={draw=gold,fill=gold!8,rounded corners=2pt,minimum width=30mm,minimum height=8mm,align=center},
  sel/.style={draw=darkgold,fill=gold!20,rounded corners=2pt,minimum width=50mm,minimum height=16mm,align=center},
  reg/.style={draw=darkgold,dashed,rounded corners=4pt,inner sep=4mm},
  arr/.style={-{Stealth[length=2.4mm]},darkgold,thick}
]
\node[box] (user) at (0,0) {Przeglądarka\\ \footnotesize sonda ping\\ \footnotesize albo ręczny wybór};
\node[box] (fe1) at (7.2,2.6) {Frontend EU};
\node[box] (be1) at (7.2,1.5) {Backend EU};
\node[box] (db1) at (7.2,0.4) {PostgreSQL EU};
\node[box] (fe2) at (7.2,-1.2) {Frontend US};
\node[box] (be2) at (7.2,-2.3) {Backend US};
\node[box] (db2) at (7.2,-3.4) {PostgreSQL US};
\begin{scope}[on background layer]
\node[reg,fit=(fe1)(be1)(db1),label={[font=\footnotesize\bfseries]above:Region EU}] (r1) {};
\node[reg,fit=(fe2)(be2)(db2),label={[font=\footnotesize\bfseries]below:Region US}] (r2) {};
\end{scope}
\draw[arr] (user.east) -- (r1.west);
\draw[arr] (user.east) -- (r2.west);
\draw[arr] (fe1) -- (be1); \draw[arr] (be1) -- (db1);
\draw[arr] (fe2) -- (be2); \draw[arr] (be2) -- (db2);
\end{tikzpicture}
\caption{Docelowa topologia. Pierwsza strona przychodzi z jednego z dwóch regionów. Skrypt na niej wybiera lobby i dalszy ruch idzie tylko tam. Trzeciego serwera nie ma.}
\end{figure}

% =====================================================================
\section{Wymagania niefunkcjonalne}
Wymagania z raportu końcowego (WN-01 do WN-12) pozostają w mocy dla \emph{każdego}
regionu. Poniżej wymagania nowe lub zmienione. Wartości docelowe są
\textbf{hipotezami} do potwierdzenia w eksperymencie F4/F5.

\begin{table}[H]
\centering\small
\begin{tabularx}{\textwidth}{@{}lYY@{}}
\toprule
\textbf{ID} & \textbf{Wymaganie} & \textbf{Miara (cel do weryfikacji)} \\
\midrule
WN-13 & Opóźnienie REST dla użytkownika w regionie własnym & p95 $<$ 200~ms (jak WN-01) \\
WN-14 & Aktualizacja listy pokojów po zdarzeniu & $<$ 1000~ms (jak WN-02), w obu regionach \\
WN-15 & Wykrycie awarii regionu i przekierowanie nowych użytkowników & Czas wykrycia i przełączenia zmierzony; cel ustalany po pierwszej próbie \\
WN-16 & Odtwarzalność środowiska & Wdrożenie z pustego konta/regionu jednym poleceniem (zrealizowane: \code{deploy/deploy.sh}, \code{deploy/azure-vm.sh}) \\
WN-17 & Spójność konfiguracji między regionami & Te same wersje obrazów z GHCR i parametryzacja w plikach \code{.run/<region>.env} \\
WN-18 & Kontrola kosztów & Budżet i alarm kosztowy; wynik kosztu w raporcie F6 \\
WN-19 & Obserwowalność & Centralny stos VictoriaLogs/Metrics i Grafana (PR 99), logi JSON z \code{request_id}, metryki kontenerów i serwer MCP \\
WN-20 & Bezpieczeństwo & Sekrety poza repozytorium, baza bez publicznego dostępu, ruch wyłącznie przez Cloudflare Tunnel, TLS/WSS \\
\bottomrule
\end{tabularx}
\end{table}

% =====================================================================
\section{Integracje}
Poniżej są tylko usługi, z których projekt korzysta. Odpowiedniki, których nie
uruchamiamy, są w sekcji~\ref{sec:chmura}.

\begin{table}[H]
\centering\small
\begin{tabularx}{\textwidth}{@{}lY@{}}
\toprule
\textbf{Element} & \textbf{Rola w rozwiązaniu} \\
\midrule
Maszyny & Trzy maszyny Azure Virtual Machines, rozmiar \code{Standard_B2ats_v2}, Ubuntu: dwie dla regionów aplikacji oraz jedna dla centralnego stosu obserwowalności \code{playlink-ops}. \\
Stos na maszynie & \code{compose.yaml}: frontend, backend, PostgreSQL i agenci (Vector, cAdvisor) w kontenerach; \code{compose.observability.yaml} dla maszyny ops. \\
Przepis regionu & Parametryzowane skrypty \code{deploy/deploy.sh} oraz konfiguracja per region w \code{.run/<region>.env} (PR 91). \\
Sekrety & \code{JWT_SECRET}, hasła bazy i \code{OBS_PASSWORD} w \code{deploy/site.env} na maszynach; brak sekretów w repozytorium. \\
Cloudflare DNS i tunel & Nazwy regionalne (\code{playlink-<reg>}), adres główny i tunel \code{playlink-ops}. Ruch wychodzący przez \code{cloudflared}; brak otwartych portów publicznych na maszynach VM. \\
Ochrona i certyfikat & Ruch wchodzi na Cloudflare (DDoS, TLS). Panel Grafany (\code{playlink-grafana.<domain>}) zabezpieczony przez Cloudflare Access (PR 99). \\
Cache plików statycznych & Cloudflare buforuje pliki JS, CSS i obrazy; sonda, API, WebSocket oraz telemetria omijają cache. \\
Wybór regionu & Sonda w przeglądarce z mechanizmem propozycji Stay/Switch (\code{RegionProposalBanner}, PR 92) oraz ręczny przełącznik. \\
Wdrożenie i aktualizacja & Publikacja wydań na GHCR (\code{release.yml}); maszyny Azure pobierają i wdrażają nowe wersje w trybie pull-based (\code{deploy/update.sh}) co 2 minuty (EU najpierw, potem US) z automatycznym rollbackiem (PR 91). \\
Obserwowalność & Grafana, VictoriaLogs (LogsQL) i VictoriaMetrics (PromQL), agenci Vector i cAdvisor (PR 99). \\
Serwer MCP & Bezpieczny serwer inspekcji i telemetryczny dla asystentów AI (\code{backend/mcp_server}) z obsługą narzędzi logów i metryk (PR 96, 97, 99). \\
Testy obciążeniowe & Dedykowany runner syntetycznego obciążenia w \code{experiments/f4/} (PR 92). \\
\bottomrule
\end{tabularx}
\end{table}

% =====================================================================
\section{Założenia i ograniczenia}

\subsection{Założenia}
\begin{itemize}
  \item Regiony: Europa (Frankfurt, spójny z domyślnym \code{eu-central}) i USA
        (wschodnie wybrzeże, spójny z \code{na-east}). \emph{Do potwierdzenia.}
  \item Pokoje są krótkotrwałe, więc ich utrata przy awarii regionu jest
        akceptowalna (RPO dla lobby bliskie zeru); trwałe dane (konta) są
        odtwarzalne z podpisu, bo tożsamość jest bezhasłowa i przechowywana wyłącznie
        po stronie użytkownika (non-custodial).
  \item Koszty pokrywają kredyty dla studentów lub niski własny budżet; eksperyment
        używa małych instancji i jest wyłączany po testach (sekcja~\ref{sec:chmura}).
  \item Frontend i backend dają się wdrożyć niezależnie (WN-10 z raportu).
\end{itemize}

\subsection{Ograniczenia}
\begin{itemize}
  \item Stan realtime w pamięci procesu (\code{ConnectionManager},
        \code{RoomChatManager}, \code{WS_LIMITS}, limiter slowapi) wymusza jedną
        instancję backendu na region albo wprowadzenie wspólnego brokera --- to
        poza zakresem.
  \item Limit REST (\code{10/minute}) jest liczony po adresie IP. Ruch wchodzi
        przez Cloudflare, więc backend czyta adres gracza z nagłówka
        \code{CF-Connecting-IP}. Bez tego wszyscy wyglądaliby jak jeden adres Cloudflare.
  \item JWT w ciasteczku \code{httpOnly} i \code{ORIGIN} muszą pasować do domeny
        regionalnej; wymaga to dopracowania konfiguracji frontendu.
  \item Nazwa użytkownika żyje w bazie regionu. Pierwsze logowanie w regionie
        dostaje losową nazwę \code{user_} i osiem znaków, tak jak dziś. Gracz
        może ją zmienić, jeśli w tym regionie jest wolna. Ta sama nazwa z drugiego
        regionu może być zajęta i wtedy zostaje przy losowej.
  \item Logowanie w drugim regionie to drugi podpis tym samym kluczem z przeglądarki.
        Token z pierwszego regionu nie przenosi sesji między osobnymi adresami.
        Wspólny \code{JWT_SECRET} zostaje, żeby token dało się zweryfikować tak samo.
  \item Czas i budżet zespołu są ograniczone terminami modułu.
\end{itemize}

\subsection{Wybór chmury i koszty}
\label{sec:chmura}
Stos aplikacji (kontenery i PostgreSQL) stawiamy na \textbf{Azure}, bo program
Azure for Students daje kredyty bez karty kredytowej i na dłuższy czas.
\textbf{Alternatywa dla stosu: AWS}, jeśli weryfikacja studencka Azure się nie
powiedzie lub kredyty okażą się niewystarczające. Konto studenckie jest już
aktywne: 88~USD do 14.10.2027. \textbf{Cloudflare} trzyma domenę, przyjmuje ruch
graczy i tunelem oddaje go maszynie. Lobby działa na maszynie, nie w Cloudflare.

\subsubsection{Dlaczego nie ma trzeciego serwera}
Dwa regiony nie wymagają maszyny, która stoi pośrodku i rozdziela ruch. Taka
maszyna byłaby wąskim gardłem: każdy gracz, także ten blisko swojego regionu,
musiałby przejść przez jeden punkt, a awaria tego punktu wyłączałaby oba regiony.

Wpisanie adresu nie łączy od razu z właściwym lobby. Nazwa jest na Cloudflare,
a tunel tej nazwy kończy się na jednej z dwóch maszyn. Stamtąd przeglądarka
dostaje stronę oraz skrypt. Oba regiony mają ten sam skrypt. Skrypt mierzy czas
odpowiedzi obu nazw i dopiero potem zmienia adres w pasku na szybszy region.
Użytkownik widzi lobby już po tym pomiarze. Ręczny przełącznik zostaje. Gdy
wybrany region przestaje odpowiadać, przeglądarka przechodzi na kolejny zdrowy
adres. Dodanie regionu to nowy adres na liście i ten sam stos z innym parametrem.
Trzeci serwer, który tylko oddawałby ten skrypt, nic nie wnosi.

Pomiar przechodzi przez Cloudflare, więc liczba milisekund jest wyższa niż przy
gołym połączeniu z maszyną. Porównanie dwóch regionów zostaje uczciwe, o ile
sonda, API i WebSocket nie są serwowane z cache. W raporcie F4 zapisujemy, że
mierzona ścieżka to przeglądarka, węzeł Cloudflare, tunel, maszyna.

\subsubsection{Co bierzemy z Cloudflare}
Wszystko poniżej jest w planie darmowym. Płatnego load balancera nie uruchamiamy.
\begin{itemize}
  \item \textbf{Tunel i DDoS.} \code{cloudflared} na każdej maszynie. Gracz nie
        łączy się z publicznym adresem maszyny. Atak na nazwę strony zatrzymuje
        się na sieci Cloudflare.
  \item \textbf{Certyfikat TLS} dla nazw regionu. Przeglądarka rozmawia po HTTPS i WSS.
  \item \textbf{Cache plików statycznych.} JS, CSS i obrazy biorą się z węzła
        blisko gracza i nie obciążają małej maszyny. Sonda, API i WebSocket
        mają regułę obejścia cache.
  \item \textbf{Kompresja i HTTP/3.} Są włączone, bo przyspieszają pobranie strony.
        Wznowienie 0-RTT zostaje wyłączone, żeby powtórzona sonda nie była
        sztucznie cichsza od pierwszej.
  \item \textbf{WebSocket.} Czat i lista pokoi przechodzą przez tunel. Istniejący
        heartbeat co 30~s zostaje, bo bezczynne połączenie Cloudflare potrafi zamknąć.
\end{itemize}
Poza zakresem zostają funkcje, które albo kosztują, albo psują pomiar i testy:
Argo Smart Routing, Load Balancing, tryb walki z botami, Rocket Loader,
kontenery Cloudflare i baza D1.

\subsubsection{Czego nie przenosimy na Cloudflare}
\begin{itemize}
  \item \textbf{Kontenery.} Cloudflare Containers nie mają planu darmowego.
        Wchodzą w Workers Paid (5~USD miesięcznie) i są rozliczane za czas
        działania. Wliczone 25~GiB-godzin nie wystarcza na dwa zawsze włączone
        lobby: nawet mały kontener działający cały miesiąc przekracza ten próg,
        a WebSocket ma pozostać otwarty. Kontenery Cloudflare są liczone od
        obsłużenia żądania do uśpienia, co nie pasuje do lobby trzymanego w
        pamięci procesu.
  \item \textbf{Baza D1.} D1 jest w planie darmowym (5~GB, limity odczytów
        i zapisów na dobę), ale jest to SQLite, nie PostgreSQL. PlayLink już
        używa PostgreSQL w obecnym stosie. Przepisanie warstwy danych
        byłoby nową aplikacją, a nie rozbudową tej, którą mamy. Hyperdrive tylko
        przyspiesza połączenie z zewnętrznym Postgresem; nie zastępuje bazy
        w regionie.
\end{itemize}

\subsubsection{Kontenery, sekrety, przepis}
Na każdy region przypada stos kontenerów: frontend, backend,
PostgreSQL oraz agenci telemetryczni (Vector, cAdvisor). Baza jest kontenerem Postgres w compose (\code{compose.yaml}). Usługi
Azure Database for PostgreSQL nie uruchamiamy. Obrazy budowane są wieloarchitekturowo i publikowane w GHCR.
Sekrety (\code{JWT_SECRET}, hasła bazy, \code{OBS_PASSWORD}) nie wchodzą do repozytorium
i nie idą do Azure Key Vault. Leżą w plikach środowiskowych na maszynie (\code{deploy/site.env}), a wdrożenie nie wymaga otwartych portów SSH dla CI.

Przepis regionu to pliki w repozytorium: \code{compose.yaml} plus plik parametrów \code{.run/<region>.env}. Jedno polecenie \code{deploy/deploy.sh eu} stawia stos regionalny, to samo polecenie z parametrem \code{us} stawia drugi region. Aktualizacje realizowane są w trybie pull-based przez timer systemd (\code{deploy/update.sh}) odpytujący rejestr GHCR, co zapewnia pełną automatyzację bez ryzyka wycieku kluczy dostępowych. Płaci się wyłącznie za czas maszyny i dysk, nie za narzędzie. Maszyna centralnego monitoringu (\code{playlink-ops}) korzysta z tego samego wzorca (\code{compose.observability.yaml}).

\textbf{Maszyna.} Usługa nazywa się Azure Virtual Machines. Rozmiar to
\code{Standard_B2ats_v2} (2 vCPU, 1~GiB RAM, AMD), obraz Ubuntu Server 24.04 LTS.
Regiony: Germany West Central i East US. Dysk systemu: 64~GB Standard SSD, bo
większy dysk Premium liczy się z kredytu osobno. Oferta obejmuje 750 godzin
miesięcznie tego rozmiaru. Jedna maszyna włączona cały miesiąc mieści się w tym
limicie, druga schodzi z kredytu 88~USD, więc po próbach obie się wyłącza. Gdy
1~GiB nie pomieści trzech kontenerów, rozmiar zmienia się na \code{Standard_B1ms}
(2~GiB) albo \code{Standard_B2s} (4~GiB). Te większe rozmiary nie są w darmowych
godzinach.

Dalsze dokumenty tego modułu mają opisywać wykonane kroki: co uruchomiono, jaki
był wynik, jaki wniosek. Sam opis dostawcy nie zamyka fazy.

\begin{table}[H]
\centering\footnotesize
\begin{tabularx}{\textwidth}{@{}lY@{}}
\toprule
\textbf{Usługa} & \textbf{Dlaczego jej nie ma w tym etapie} \\
\midrule
Azure Database for PostgreSQL & Baza już jest kontenerem Postgres w compose. Osobna usługa zjadłaby kredyt i zdublowała dane. \\
Azure Key Vault & Sekrety już obsługuje GitHub Actions i plik na maszynie. \\
Azure Traffic Manager & Wybór regionu robi skrypt i zmienia adres w pasku. Traffic Manager tylko podmienia odpowiedź DNS, nie chroni maszyny i potrafi przerzucić gracza na drugi region pod tą samą nazwą. \\
Cloudflare Load Balancing i Argo & Płatne. Tunel i skrypt robią ochronę oraz wybór regionu bez abonamentu. \\
\bottomrule
\end{tabularx}
\caption{Usługi rozważane i świadomie pominięte.}
\label{tab:mapa}
\end{table}

\begin{table}[H]
\centering\footnotesize
\begin{tabularx}{\textwidth}{@{}lYY@{}}
\toprule
 & \textbf{Azure for Students} & \textbf{AWS} \\
\midrule
Kredyty & 100~USD na 12 miesięcy, odnawiane co roku dla studenta & Do 200~USD: 100~USD na start oraz po 20~USD za pięć zadań (uruchomienie instancji EC2, baza RDS, funkcja Lambda, zapytanie w Bedrock, ustawienie budżetu); plan darmowy wygasa po 6 miesiącach lub po wyczerpaniu kredytów \\
Karta kredytowa & Nie jest wymagana & Wymagana przy zakładaniu konta \\
Weryfikacja studenta & Poczta uczelniana; domena musi być rozpoznana & Nie dotyczy (kredyty dla nowych kont); AWS Educate osobno \\
Ryzyko & Uczelniana domena może nie zostać rozpoznana & Po wyczerpaniu kredytów płatne zasoby; wymagany budżet z alarmem \\
\bottomrule
\end{tabularx}
\caption{Kredyty dla studentów (dane orientacyjne z publicznych opisów programów; przed założeniem konta zweryfikować na stronach dostawców).}
\end{table}

\textbf{Szacunek.} Z kredytu schodzi czas drugiej maszyny i dyski. Jedna maszyna
tego rozmiaru mieści się w 750 darmowych godzinach miesięcznie. Baza w kontenerze
osobnego rachunku nie generuje. Cloudflare w zakresie powyżej jest w planie
darmowym. Po próbach maszyny się wyłącza. Alarm budżetu ustawiamy, zanim druga
maszyna zostanie włączona na dłużej. F6 wpisze kwotę z rzeczywistego rachunku.

\textbf{Decyzja fazy 0} jest zamknięta: Azure for Students, maszyny
\code{Standard_B2ats_v2}, baza w kontenerze, Cloudflare przy nazwach i tunelu.
AWS zostaje tylko awaryjnie, gdyby kredyt się skończył przed końcem modułu.

% =====================================================================
\section{Ryzyka}
\begin{table}[H]
\centering\footnotesize
\begin{tabularx}{\textwidth}{@{}lYcY@{}}
\toprule
\textbf{ID} & \textbf{Ryzyko} & \textbf{Wpływ} & \textbf{Reakcja} \\
\midrule
R1 & Podział społeczności: gracze z różnych regionów nie widzą wzajemnych lobby & Wysoki & Ręczny wybór regionu; podgląd listy drugiego regionu jako rozszerzenie; ujęcie w F6 \\
R2 & Koszty przekraczają budżet lub kredyty (transfer, baza, kontrola stanu) & Średni & Małe instancje, budżet z alarmem, wyłączanie środowiska po testach; zapasowo drugi dostawca (AWS) \\
R3 & Stan realtime w pamięci uniemożliwia skalowanie w obrębie regionu & Średni & Jedna instancja na region w zakresie; udokumentowane jako ograniczenie \\
R4 & Różnice konfiguracji między regionami zafałszowują eksperyment & Wysoki & Jeden szablon, te same obrazy, porównanie konfiguracji przed testem \\
R5 & Zła metodologia pomiarów (generator blisko serwera, krótki test, szum) & Wysoki & Generatory w obu kontynentach, powtórzenia, opis warunków \\
R6 & Zbyt wolne wykrycie awarii przy domyślnych ustawieniach kontroli stanu & Średni & Zmierzyć i dostroić; wynik zapisać w F5 \\
R7 & Dane użytkowników w USA a wymagania prawne (RODO) & Średni & Przechowywane są tylko adresy tożsamości i nazwy; ocena w F6, nie wdrożenie zgodności \\
R8 & Zakres ponad możliwości zespołu w 10 tygodniach & Średni & Priorytety: F1--F3 w pierwszej kolejności, F6 zawsze dostarczone \\
\bottomrule
\end{tabularx}
\end{table}

% =====================================================================
\section{Kryteria sukcesu}
\begin{table}[H]
\centering\small
\begin{tabularx}{\textwidth}{@{}lYY@{}}
\toprule
\textbf{Filar} & \textbf{Kryterium sukcesu} & \textbf{Dowód i stan realizacji} \\
\midrule
F1 & Oba regiony przechodzą ten sam zestaw testów przepływu (pokój, dołączenie, czat, lista) & Spełnione: stosy EU/US za Cloudflare Tunnel, testy backendu i frontendu przechodzą (PR 91) \\
F2 & Pierwsze wejście trafia do szybszego regionu; szybszy drugi region jest propozycją (Stay/Switch) & Spełnione: baner \code{RegionProposalBanner}, cookie \code{playlink.region}, failover przy awarii (PR 91, 92) \\
F3 & Czyste wdrożenie obu środowisk jednym poleceniem; powtarzalna aktualizacja wydań & Spełnione: \code{deploy/deploy.sh}, obrazy GHCR i timer pull-based \code{deploy/update.sh} (PR 91) \\
F4 & Zebrane pomiary wariantów A i B dla obu grup użytkowników, z opisem warunków & Spełnione: moduł \code{experiments/f4/}, raporty wyników p50/p95, błędy i WebSocket (PR 92) \\
F5 & Co najmniej dwa scenariusze awarii wykonane, z czasem wykrycia i odzyskania & Scenariusze stopu backendu/tunelu i procedury runbookowe (\code{docs/operations/runbook-regions.md}) \\
F6 & Raport z rekomendacją: kiedy opłaca się drugi region, jaki koszt i jakie ograniczenia & Podsumowanie wyników eksperymentu F4, kosztów maszyn Azure oraz wniosków \\
F7 & Centralna obserwowalność oraz bezpieczna inspekcja stanu przez asystentów AI (MCP) & Spełnione: Grafana, VictoriaLogs/Metrics (PR 99), serwer MCP z narzędziami logów i metryk (PR 96, 97) \\
\bottomrule
\end{tabularx}
\end{table}
\textbf{Hipotezy do sprawdzenia:} H1 --- dla użytkowników z USA wariant A (tylko Europa)
nie spełnia WN-13, a wariant B tak. H2 --- dla użytkowników z Europy wariant B nie
pogarsza wyników względem A. H3 --- po awarii jednego regionu drugi pozostaje w
pełni użyteczny. Wynik negatywny też jest wartościowym wnioskiem F6.

% =====================================================================
\section{Roadmapa produktu}
Plan obejmuje 10 tygodni liczonych od startu prac. Architektura na ten etap jest
ustalona. Konto studenckie jest aktywne (88~USD, ważność do 14.10.2027).
Stos staje na maszynach wirtualnych \code{Standard_B2ats_v2}. Reszta tygodni jest wykonaniem, pomiarem
i zapisem wyniku.

\subsection{Dziennik etapu}
Każdy tydzień zamyka krótki wpis, ten sam dla wszystkich faz. Wpis ma cztery pola:
co było zaplanowane, co zostało uruchomione, jaki był wynik (liczba, czas, zrzut
albo komunikat błędu) i jaki z tego wniosek, także gdy coś nie działa. Sam opis
dostawcy nie zamyka tygodnia. Z tych wpisów składa się materiał dla prowadzącego
oraz raport F6 (prowadzony na bieżąco w \code{docs/dziennik-etapu.md}).

\begin{table}[H]
\centering\footnotesize
\begin{tabularx}{\textwidth}{@{}llYY@{}}
\toprule
\textbf{Tydzień} & \textbf{Filar} & \textbf{Praca zrealizowana} & \textbf{Wpis w dzienniku i rezultaty} \\
\midrule
1 & decyzje & Konto Azure for Students aktywne; maszyny \code{Standard_B2ats_v2} z Ubuntu; konfiguracja domen i tuneli Cloudflare & Zrzut salda studenckiego, specyfikacja maszyn, domeny w Cloudflare \\
2--4 & F1, F3 & Wdrożenie wieloregionowe EU/US za tunelami Cloudflare; potok budowania wydań GHCR i aktualizacja pull-based (\code{deploy/update.sh}) & Działające dwa regiony, healthchecki, runbook operacyjny (\code{docs/operations/runbook-regions.md}) (PR 91) \\
5--7 & F2, F4 & Lokalne środowisko wieloregionowe (\code{start-local.ps1}), propozycja Stay/Switch (F2) w UI oraz realizacja eksperymentu obciążeniowego F4 w \code{experiments/f4/} & Pomiary p50/p95 wariantu 1 vs 2 regionów, weryfikacja braku samoczynnego przeskakiwania lobby (PR 92) \\
8 & MCP & Implementacja serwera Model Context Protocol: bezpieczna inspekcja gier i pokoi, dynamiczne regiony RFC 6570 (\code{backend/mcp_server}) & Narzędzia i zasoby MCP przetestowane lokalnie i w CI (PR 96, PR 97) \\
9 & Obserw. & Centralny stos obserwowalności na maszynie \code{playlink-ops} (VictoriaLogs, VictoriaMetrics, Grafana) oraz narzędzia logów i metryk w MCP & Dashboard \emph{PlayLink regions}, logi z \code{request_id}, narzędzia \code{search_logs} i \code{query_metrics} w MCP (PR 99) \\
10 & F5, F6 & Próby awarii, podsumowanie kosztów rzeczywistych, raport końcowy i odbiór projektu & Komplet wyników w dzienniku etapu, wnioski architektoniczne i obrona projektu \\
\bottomrule
\end{tabularx}
\end{table}
Poza horyzontem tego etapu: podgląd lobby drugiego regionu,
wspólny broker zdarzeń i wdrożenie trzeciego regionu z gotowego szablonu.

% =====================================================================
\section{Korzyści biznesowe}
\begin{itemize}
  \item \textbf{Dla graczy:} niższe opóźnienia i dostęp do lobby blisko siebie lub
        w drugim regionie na życzenie.
  \item \textbf{Dla ciągłości usługi:} awaria jednego regionu nie wyłącza całej
        platformy.
  \item \textbf{Dla zespołu:} powtarzalne wdrożenie zamiast ręcznej konfiguracji
        jednego hosta; łatwiejsze dodanie kolejnego środowiska.
  \item \textbf{Dla decyzji:} dane, czy drugi region jest opłacalny, zanim zostaną
        poniesione stałe koszty.
  \item \textbf{Dla oceny modułu:} jasno rozdzielone, mierzalne rezultaty oraz
        nabyta praktyka z chmury, IaC i testowania odporności.
\end{itemize}

% =====================================================================
\section{Załączniki}
\begin{enumerate}[label=Z\arabic*.]
  \item Raport końcowy PlayLink (\code{prezentacje/raport_koncowy/raport.pdf}).
  \item Podręcznik wdrożenia i operacji wieloregionowych: \\ \code{docs/operations/runbook-regions.md}.
  \item Dokumentacja obserwowalności (Grafana, VictoriaLogs, VictoriaMetrics): \\ \code{docs/operations/observability.md}.
  \item Dokumentacja serwera Model Context Protocol (MCP): \\ \code{backend/mcp_server/README.md}.
  \item Moduł i raport eksperymentu obciążeniowego 1 vs 2 regiony: \code{experiments/f4/}.
  \item Dziennik postępów etapu: \code{docs/dziennik-etapu.md}.
  \item Potok wydań i publikacji obrazów: \code{.github/workflows/release.yml}.
  \item Lista lokalizacji lobby: \code{backend/config.py} (\code{LOBBY_LOCATIONS}).
\end{enumerate}

\end{document}
